\documentclass[a4paper,12pt]{article}
\usepackage{amsmath,amssymb,amsthm}
\usepackage[margin=1in]{geometry}

\newtheorem{theorem}{Theorem}
\newtheorem{lemma}[theorem]{Lemma}

\title{Problem 347: Largest Integer Divisible by Two Primes}
\author{Project Euler}
\date{}

\begin{document}
\maketitle

\section{Problem Statement}
For distinct primes $p < q$, let $M(p,q,N)$ be the largest integer $\le N$ whose only prime factors are $p$ and $q$ (both appearing). Find
\[
\sum_{\substack{p < q \\ pq \le 10^7}} M(p, q, 10^7).
\]

\section{Mathematical Foundation}

\begin{theorem}[Structure]
The integers with only prime factors $p$ and $q$ (both present) form the set $\{p^a q^b : a \ge 1,\, b \ge 1\}$.
\end{theorem}

\begin{proof}
By unique factorization, an integer whose only prime factors are $p$ and $q$ has the form $p^a q^b$ with $a, b \ge 0$. Requiring both primes to appear forces $a \ge 1$ and $b \ge 1$.
\end{proof}

\begin{lemma}[Finite Enumeration]
For $pq \le N$, the set $\{p^a q^b \le N : a,b \ge 1\}$ has at most $\log_p N \cdot \log_q N = O(\log^2 N)$ elements.
\end{lemma}

\begin{proof}
$a \le \log_p(N/q) \le \log_p N$ and $b \le \log_q(N/p) \le \log_q N$.
\end{proof}

\begin{theorem}[Maximum Exists]
$M(p,q,N) = \max\{p^a q^b : a,b \ge 1,\; p^a q^b \le N\}$ is well-defined when $pq \le N$.
\end{theorem}

\begin{proof}
The set is non-empty ($pq \le N$) and finite ($p^a q^b \ge 2^{a+b} \to \infty$). A finite non-empty subset of $\mathbb{Z}$ has a maximum.
\end{proof}

\begin{lemma}[Pair Count]
The number of prime pairs $(p,q)$ with $p < q$ and $pq \le N$ is $O(N/\log^2 N)$.
\end{lemma}

\begin{proof}
For each prime $p \le \sqrt{N}$, the primes $q \in (p, N/p]$ number $O(N/(p \log N))$ by PNT. Summing over $p$ gives $O(N/\log^2 N)$.
\end{proof}

\section{Algorithm}

\begin{verbatim}
primes = sieve(N/2)
total = 0
for each pair (p, q) with p < q, p*q <= N:
    best = 0
    pp = p
    while pp * q <= N:
        pq = pp * q
        while pq * q <= N: pq *= q
        best = max(best, pq)
        pp *= p
    total += best
return total
\end{verbatim}

\section{Complexity Analysis}
\begin{itemize}
    \item \textbf{Time:} $O(N)$. Each of $O(N/\log^2 N)$ pairs requires $O(\log^2 N)$ work.
    \item \textbf{Space:} $O(N)$ for the sieve.
\end{itemize}

\section{Answer}

\[
\boxed{11109800204052}
\]

\end{document}
